% !TeX root = ../papers/nc-p3.tex
\section{内側衝突の残りの錐と単元座標の極}
\label{sec:rank-fifteen-inner-collision-fan}
\begin{theorem}\label{thm:rank-fifteen-inner-collision-fan}
標数零の通常な階数15窓が剰余標数3において $\nu(x)=0$, $\nu(y)>0$ を満たせば、
標準Koszul左右AS正則で、各始点で四中心四次式上自由階数256である。
\end{theorem}
\begin{proof}
第\ref{sec:rank-fifteen-unit-pole-depth}節と第\ref{sec:rank-fifteen-outer-pole-fan}節により、
$\bar x=1$, $\bar J=1$ だけを扱えばよい。前節の $z=1+\delta$ の表示を使い、
$(\nu(H),\nu(\delta))$ の順に傾きを書く。
$(1,2)$ の半直線とそれより深い衝突、および $(1,1)$ の半直線は前節で決まった。
残る範囲は境界 $(1,0)$ から $(2,1),(1,1),(1,2)$ までの三つの開いた錐と
$(2,1)$ の半直線に尽きる。

$(2,1)$ では $H=m\delta^2$、$m$ は単元である。$z=1+\delta$ のまま
$\mathbb Q(m)$ 上で三十桁を使うと、長さ1から5の損失は $(0,4,4,5,4)$、
四次中心は $(0,1,1,1)$、五次の積は一桁を使う。
長さ6から13では $(3,0,2,1,2,0,2,0)$ 桁を使い、
主元の定数分子・分母の剰余の平方自由な積は $m(m+1)$ となる。
$m+1$ の上で単元の主元に取り直した同じ標数零計算ではその積は $m(m-1)$ となる。
二表示で全単元を覆う。

三つの錐で両辺の原始ベクトルの行列式を取るといずれも1であり、外側と同じ単項式の
座標変換が使える。対角線は順に $(3,1),(3,2),(2,3)$ である。
\path{computations/verify_rank_fifteen_inner_collision_fan.py} は
これらの対角線でそれぞれ46、44、30桁を用いて全308組の指数と中心座標を求める。
$(1,0)$ の辺では $z$ が一般の接線座標、$(2,1)$ では $H/\delta^2$ が一般元、
$(1,1)$ と $(1,2)$ では $\delta/H$ と $\delta/H^2$ がそれぞれ一般元である。

今回の錐では、同じ文字の二つの中心式を選ぶ際の座標が辺によって変わる。
例えば最後の錐の対角線では、文字0の三座標のうち一つは始点2、残る二つは始点3にある。
一方 $\delta=mH$ のもとの表示では三座標とも始点3にある。
このため、各錐で対角線上の中心座標を使う。
座標を指定する各四次語基底を辺の既知の経路格子に写し、単元主元で
正方行列が可逆であることを確認する。この基底で中心方程式を作り直し、
指定列を自由座標として残すように他の列から主元を取る。
その後の四中心商はこの基底で指定した同じ二本ずつの中心式から計算する。
境界 $(1,0)$ は最初の対角線の座標で、辺 $(1,1)$ は最後の錐に入る際に
その対角線の座標でも、標数零の有理関数体上で検証する。

これにより各錐の両辺と対角線で選んだ中心は同じ規格化を持つ。
各行列について昇順に並べた両辺の指数の和が対角線の全指数と一致する。
未切断経路から四次中心、五次までの中心商、左右の積表の順に
補題\ref{lem:rank-fifteen-two-parameter-smith}を適用すれば、
所定の正付値の両座標で十三次の消滅までの上界が得られる。
従来の全次数判定から結論が従う。
\end{proof}

\begin{corollary}
$x,y\in\mathbb Q$, $\nu_3(x)=0$, $y\ne0$ の通常な階数15窓は、
全て標準Koszul左右AS正則で、各始点で四中心四次式上自由階数256である。
\end{corollary}
\begin{proof}
$\nu_3(y)>0$ は本節、$\nu_3(y)=0$ は素体単元の定理、
$\nu_3(y)<0$ は $\nu_3(y)<\min(0,2\nu_3(x))$ の極の定理である。
\end{proof}

\begin{corollary}\label{cor:rank-fifteen-unit-vertical-conics}
3の上のある付値で $\bar a=1$ とする。標数零の代数閉体上の縦切片 $x=a$ では、
$y\ne0$ の全ての水準が標準Koszul左右AS正則である。
$f=(1-a)^2-4a^3/27$ と置くと、添字保持同型類は
\[
      T\in\mathbb G_m,\qquad
      y=\frac{a^3}{27}-\frac{(T-f/T)^2}{16}\ne0
\]
で重複なく表される。非零水準のうち $a^3/27$ または $(1-a)^2/4$ に等しいものは二類、
他は四類であり、全類の最小周期は2である。特に $a\in\mathbb Q$, $a\equiv1\pmod3$
なら $y$ は任意の非零な体の元でよい。
\end{corollary}
\begin{proof}
付値を係数体へ延長する。正付値の $y$ は本節の定理、
単元は定理\ref{thm:rank-fifteen-vertical-all-levels}、負付値は極の定理が覆う。
$\bar a=1$ なので $a\ne3,3/4$、従って $f\ne0$ である。
定理\ref{thm:rank-fifteen-formal-boundary}の残り二点の横座標は
$4a^2-12a+7=0$ を満たし、$\bar a=1$ とは両立しない。
従ってこの縦切片の全ての点は通常表示をもつ。$p^2=a^3-27y$,
$q^2=(1-a)^2-4y$ と全点での三不変量の分離を用いると、
$T=q+2p/(3\sqrt3)$, $f/T=q-2p/(3\sqrt3)$ から上の式を得る。
$f\ne0$ により二つの分岐水準は異なる。$p,q$ の一方の平方が零なら二類、
両方非零なら四類となる。通常窓の最小周期は2である。
\end{proof}
